% =====================================================================
%  Fig. 1 -- the flowchart alphabet, defined once for the whole report.
%  Spans both columns (figure*) because it is a reference the reader
%  comes back to.
% =====================================================================
\begin{figure*}[!t]
\centering
\begin{tikzpicture}[fcchain]

% ---------------------------------------------------------------------
%  LEFT BLOCK : the six node shapes
% ---------------------------------------------------------------------
\node[fcterm] (s1) {START / STOP};
\node[fcproc, below=of s1] (s2) {Process\\{\tiny a computation}};
\node[fcdec,  below=of s2] (s3) {Decision?};
\node[fcio,   below=8mm of s3] (s4) {ROS topic in / out};
\node[fcsub,  below=of s4] (s5) {lib/ call\\{\tiny framework-independent}};
\node[fcfail, below=of s5] (s6) {Recovery / degraded\\{\tiny branch taken on failure}};

\node[font=\sffamily\footnotesize\bfseries, above=3mm of s1] (shapeshdr)
     {(a) Shapes};

% shape annotations, to the right of each
\node[anchor=west, font=\sffamily\scriptsize, right=3mm of s1]
     {Entry and exit points of the node.};
\node[anchor=west, font=\sffamily\scriptsize, right=3mm of s2]
     {Ordinary computation, no I/O.};
\node[anchor=west, font=\sffamily\scriptsize, right=3mm of s3]
     {Test; both exits are labelled \textsf{yes}/\textsf{no}.};
\node[anchor=west, font=\sffamily\scriptsize, right=3mm of s4]
     {A message published or received.};
\node[anchor=west, font=\sffamily\scriptsize, right=3mm of s5]
     {Call into \filename{lib/}: no ROS, no simulator.};
\node[anchor=west, font=\sffamily\scriptsize, right=3mm of s6]
     {Fallback path: brake, escape, abandon goal.};

% ---------------------------------------------------------------------
%  RIGHT BLOCK : the three phase bands
% ---------------------------------------------------------------------
\node[fcterm, right=78mm of s1] (p0) {node start};
\node[fcproc, below=of p0] (p1) {declare parameters,\\open publishers \& subscribers};
\node[fcproc, below=of p1] (p2) {arm timers, seed state};

\node[fcio,   below=9mm of p2] (p3) {wait for next message / tick};
\node[fcproc, below=of p3] (p4) {read inputs, run the algorithm};
\node[fcdec,  below=of p4] (p5) {shutdown\\requested?};
\node[fcio,   below=9mm of p5] (p6) {publish command};

\node[fcproc, below=10mm of p6] (p7) {stop the base, release the arm};
\node[fcproc, below=of p7] (p8) {emit final report};
\node[fcterm, below=of p8] (p9) {destroy node, rclpy.shutdown()};

\draw[fcline] (p0) -- (p1);
\draw[fcline] (p1) -- (p2);
\draw[fcline] (p2) -- (p3);
\draw[fcline] (p3) -- (p4);
\draw[fcline] (p4) -- (p5);
\draw[fcline] (p5) -- node[fclab,right] {no} (p6);
\draw[fcback] (p6.east) -- ++(14mm,0) |- (p3.east);
\node[font=\sffamily\tiny\itshape, text=phaseMainLine]
     at ($(p6.east)+(20mm,7mm)$) {loop};
\draw[fcline] (p5.west) -- ++(-13mm,0) |- node[fclab,pos=0.25] {yes} (p7.west);

\node[font=\sffamily\footnotesize\bfseries, above=3mm of p0] {(b) Phase bands};

% ---- the bands themselves, on the background layer -------------------
\begin{scope}[on background layer]
  \phaseband{bandinit}{INITIALISATION}{phaseInitLine}{(p0)(p1)(p2)}{bi}
  \phaseband{bandmain}{MAIN LOOP}{phaseMainLine}{(p3)(p4)(p5)(p6)}{bm}
  \phaseband{bandend}{TERMINATION / DESTRUCTION}{phaseEndLine}{(p7)(p8)(p9)}{be}
\end{scope}

% ---- band annotations ------------------------------------------------
\node[fcnote, right=6mm of bi.east, anchor=west, text width=44mm]
     {\textbf{Executed once.} Everything whose value must survive the
      whole run. Parameters are declared here so a YAML file can
      override them --- the source of one of the regressions reported in
      Sec.~\ref{sec:regression}.};
\node[fcnote, right=6mm of bm.east, anchor=west, text width=44mm]
     {\textbf{Executed at a fixed rate.} One tick, one decision. State
      that must be re-armed per cycle belongs here and not above: a
      counter armed once in initialisation is exactly the defect that
      disabled multi-pick harvesting (Sec.~\ref{sec:phaseB_mission}).};
\node[fcnote, right=6mm of be.east, anchor=west, text width=44mm]
     {\textbf{Executed once, on the way out.} Must run even on
      \texttt{SIGINT}: the performance monitor prints its final report
      from this band, because every run of this project ends with
      Ctrl\,+\,C.};

\end{tikzpicture}
\caption{The flowchart language used throughout this report, defined
once. (a) The six node shapes: two are specific to a robotics stack ---
a \emph{ROS topic} shape, so the boundary with the middleware is
visible, and a \emph{library call} shape, so the boundary with the
framework-independent algorithm layer is visible. (b) The three phase
bands into which every later flowchart is partitioned. The banding is
functional rather than decorative: two of the defects analysed in this
report are cases of code placed in the wrong band.}
\label{fig:legend}
\end{figure*}
